% =============================================================================
%  text/THE ARGUMENT/part-1.tex
%  Demonstrates theorem-like environments (assumption, proposition,
%  hypothesis). The statements below are generic math/logic examples and are
%  unrelated to any research topic -- replace them with your own statements.
%  演示假设、命题、假设（hypothesis）等环境；以下为通用数学示例，请替换。
% =============================================================================

This subsection demonstrates the theorem-like environments provided by the
template: \texttt{theorem}, \texttt{lemma}, \texttt{proposition},
\texttt{corollary}, \texttt{assumption}, and \texttt{hypothesis}. The
examples below are intentionally generic.

\begin{assumption}
All sets considered below are finite.
\end{assumption}

\begin{proposition}
For any two sets $A$ and $B$,
\[
|A \cup B| = |A| + |B| - |A \cap B|.
\]
\end{proposition}

\begin{hypothesis}[Example hypothesis]
The two statements below are logically equivalent:
\[
(A \subseteq B) \iff (A \cap B = A).
\]
\label{hyp:example}
\end{hypothesis}

Hypotheses can be numbered and referenced, as in Hypothesis~\ref{hyp:example}.
